\documentclass{jlreq}
\usepackage[utf8]{inputenc}
\usepackage[T1]{fontenc}
\usepackage{geometry}
\usepackage{graphicx}
\usepackage{float}
\usepackage{hyperref}
\usepackage[english,indonesian]{babel}

\title{Optimasi Pemrosesan Template LaTeX Menggunakan Pendekatan Berbasis AI}
\author{Academic Support Agent}
\date{\today}

\begin{document}
\maketitle

\begin{abstract}

\end{abstract}

\chapter{PENDAHULUAN}


\chapter{TINJAUAN PUSTAKA}


% Original content (commented out):
% \documentclass{jlreq}

\usepackage{hyperref}

\title{日本語 jlreq class  基本テンプレート}
\author{東京工業大学 情報理工学院}
\date{May 2021}

\begin{document}

\maketitle

\section{Overleaf での使い方}
Compiler は Lua\LaTeX を選択してください．

jlreq は \TeX Live でプリインストールされているパッケージなので、通常通り

\verb+\documentclass{jlreq}+

とするのみでパッケージを呼び込めます．これで横書きのarticle相当の文書クラスとなります．縦書きにするには tate オプションを渡します．また， report や book 相当の文書クラスとするには，それぞれ report や book オプションを渡します．たとえば，縦書きの本を作成するには

\verb+\documentclass[tate,book]{jlreq}+

とします．

その他， \verb+oneside / twoside / onecolumn / twocolumn / titlepage / notitlepage / draft /+ \\ \verb+final / openright / openany / leqno / fleqn+ というよくあるオプションを受け付けます．

標準的な文書クラスと同じように中身を書くことができます．

\section{公式サイト}
\begin{itemize}
    \item \href{https://github.com/abenori/jlreq}{Github -- abenori/jlreq}
    \item \href{https://www.tug.org/texlive//Contents/live/texmf-dist/doc/latex/jlreq/jlreq-ja.html}{TUG -- jlreq} 本テンプレートの使い方の説明はこのページに基づく
\end{itemize}
\end{document}


\end{document}